\textbf{COLUNI 2025 -- Questão 05}\par\medskip Leia o trecho a seguir e faça o que se pede.

\par

Belonísia me encontrou dobrando algumas peças de roupa e guardando na mala que pertenceu a nossa avó. Vi seus olhos surpresos com a descoberta e não fui capaz de comunicar nada sobre o que estava fazendo. Nem ela. Seu olhar era inquisidor, árido como o tempo que nos cercava, e minha vergonha era suficiente para que sequer tentasse justificar o que fazia. Chorei depois que minha irmã saiu, porque tinha a certeza de que estava contribuindo para seu sofrimento. Estava tendo encontros com nosso primo, talvez ela já desconfiasse da gravidez, mas além de tudo isso, eu era sua irmã – não tínhamos segredos, ou ao menos evitávamos ter – e havia me encerrado em meu mundo naquelas últimas semanas, esquecendo-me da família e dela, principalmente, que parecia cada vez mais distante. (VIEIRA JUNIOR, Itamar. \textit{Torto arado}. São Paulo: Todavia, 2019. p. 82)

\par

Com base na leitura do trecho acima, é CORRETO afirmar que retrata:\begin{enumerate}
\item a felicidade de Bibiana e sua indiferença em relação à irmã.
\item a vergonha de Bibiana e seu contentamento em relação à partida.
\item a realização do sonho de Bibiana e o desejo de retornar à infância.
\item a tristeza de Bibiana e seu grande amor por Belonísia.\end{enumerate}
